% =============================================================
% CONFIGURAÇÃO GERAL - UFSC / ABNT
% Compilar com pdfLaTeX + BibTeX (ou latexmk -pdf).
% =============================================================

\usepackage[T1]{fontenc}
\usepackage[utf8]{inputenc}
% O template Word utiliza Arial. Helvetica (pacote helvet) é a aproximação
% livre e portátil adotada aqui para compilar sem arquivos de fonte externos.
\usepackage[scaled=0.97]{helvet}
\renewcommand{\familydefault}{\sfdefault}

\usepackage{microtype}
\usepackage{graphicx}
\usepackage{booktabs}
\usepackage{longtable}
\usepackage{tabularx}
\usepackage{array}
\usepackage{multirow}
\usepackage{amsmath,amssymb,mathtools,bm}
\usepackage{siunitx}
\usepackage{float}
\usepackage{enumitem}
\usepackage{caption}
\usepackage{subcaption}
\usepackage{xcolor}
\usepackage{url}
\usepackage{hyperref}
\usepackage[alf]{abntex2cite}

% -------------------------------------------------------------
% Página e parágrafos - conforme o template fornecido
% Superior/esquerda: 3 cm; inferior/direita: 2 cm.
% Corpo: 12 pt, espaçamento 1,5; recuo de primeira linha: 1,5 cm.
% -------------------------------------------------------------
\setlrmarginsandblock{3cm}{2cm}{*}
\setulmarginsandblock{3cm}{2cm}{*}
\checkandfixthelayout

\setlength{\parindent}{1.5cm}
\setlength{\parskip}{0pt}
\OnehalfSpacing

% -------------------------------------------------------------
% Títulos
% -------------------------------------------------------------
\renewcommand{\ABNTEXchapterfont}{\bfseries\normalsize}
\renewcommand{\ABNTEXchapterfontsize}{\normalsize}
\renewcommand{\ABNTEXsectionfont}{\normalfont\normalsize}
\renewcommand{\ABNTEXsectionfontsize}{\normalsize}
\renewcommand{\ABNTEXsubsectionfont}{\normalfont\normalsize}
\renewcommand{\ABNTEXsubsectionfontsize}{\normalsize}
\renewcommand{\ABNTEXsubsubsectionfont}{\normalfont\normalsize}
\renewcommand{\ABNTEXsubsubsectionfontsize}{\normalsize}

% Capítulos iniciam em nova página (comportamento padrão da classe).
\setbeforesecskip{1.5\onelineskip}
\setaftersecskip{0.5\onelineskip}


% -------------------------------------------------------------
% Aparência dos capítulos e dos títulos pré-textuais.
%
% Capítulos numerados: alinhados à esquerda, no formato "1 - INTRODUÇÃO".
% Elementos pré-textuais não numerados: centralizados, em negrito e caixa alta,
% conforme o template da Engenharia Mecânica/UFSC.
% -------------------------------------------------------------
\renewcommand{\printchaptername}{}
\renewcommand{\chapternamenum}{}
\renewcommand{\printchapternum}{\thechapter.\space}
\renewcommand{\afterchapternum}{}

\makeatletter
\renewcommand{\printchaptertitle}[1]{%
  \ifthenelse{\boolean{abntex@innonumchapter}}{%
    \centering\bfseries\normalsize\MakeUppercase{#1}%
  }{%
    \raggedright\bfseries\normalsize\MakeUppercase{#1}%
  }%
}
\makeatother

% O ambiente de resumo do abnTeX2 usa uma formatação própria para o título.
% A definição abaixo garante o mesmo padrão: negrito, centralizado e caixa alta.
\renewcommand{\abstractnamefont}{\bfseries\normalsize}

% -------------------------------------------------------------
% Legendas
% Figuras: legenda abaixo (controle pela posição de \caption).
% Tabelas: legenda acima (controle pela posição de \caption).
% -------------------------------------------------------------
\captionsetup{font=small,labelfont=normalfont,labelsep=endash,justification=centering}

% -------------------------------------------------------------
% Numeração no canto superior direito apenas a partir da Introdução.
% -------------------------------------------------------------
\makepagestyle{ufsc}
\makeoddhead{ufsc}{}{}{\normalsize\thepage}
\makeevenhead{ufsc}{}{}{\normalsize\thepage}
\makeoddfoot{ufsc}{}{}{}
\makeevenfoot{ufsc}{}{}{}
\aliaspagestyle{chapter}{ufsc}

% -------------------------------------------------------------
% Links discretos no PDF.
% -------------------------------------------------------------
\hypersetup{
    pdftitle={\tituloTCC},
    pdfauthor={\autorTCC},
    colorlinks=true,
    linkcolor=black,
    citecolor=black,
    urlcolor=black,
    bookmarksdepth=4
}

% -------------------------------------------------------------
% Atalhos úteis
% -------------------------------------------------------------
\newcommand{\propeller}{\textit{propeller}}
\newcommand{\propellers}{\textit{propellers}}
\newcommand{\todo}[1]{\textbf{[TODO: #1]}}
